\documentclass{plantillaPPT}

\titulo{2}{Integral Indefinida (Parte II)}

\begin{document}

\primeraDiapo

\begin{frame}{Verdadero o Falso}
    \framesubtitle{Práctica 2}
    \vspace{-1cm}
1. Decida si las siguientes afirmaciones son verdaderas o falsas. Justifique su respuesta. \vspace{1cm}

\((a) \; \; \text{Si } f: \mathbb{R} \to \mathbb{R}\) es una función derivable y positiva, entonces

\[\int \frac{f'(1+2x)}{f(1+2x)}dx = 2\ln(f(1+2x)) + C\]

    
\end{frame}

\begin{frame}{Verdadero o Falso}
    \framesubtitle{Práctica 2}

\((b)\) Usando integración por partes, con:
\[u=\frac{1}{x} \Rightarrow du= \frac{-1}{x^2}dx \qquad y \qquad 
dv = dx \Rightarrow v=x\]

se obtiene \(\int \frac{1}{x}dx = 1 + \int \frac{1}{x}dx\), cancelando la integral en cada miembro resulta \(1=0\)
    
\end{frame}

\begin{frame}{Integral Indefinida}
    \framesubtitle{Práctica 2}
    \vspace{-1cm}

2. Calcule las siguientes integrales indefinidas:
\vspace{0.5cm}

\[(a) \int e^{\arccos(x)}dx\]

\end{frame}

\begin{frame}{Integral Indefinida}
    \framesubtitle{Práctica 2}

\[(b) \int \sqrt{2x-x^2}dx\]
    
\end{frame}

\begin{frame}{Integral Indefinida}
    \framesubtitle{Práctica 2}

\[(c) \int \arctan(\sqrt{x}) dx\]
    
\end{frame}

\end{document}